% =============================================================
%  CAHIER DES CHARGES — SAP Performance Dashboard
%  Version fonctionnelle — orientée métier et utilisateurs
%  Compilable sur Overleaf (pdflatex) — UTF-8
% =============================================================
\documentclass[12pt,a4paper]{article}

% ── Encodage & langue ─────────────────────────────────────────
\usepackage[utf8]{inputenc}
\usepackage[T1]{fontenc}
\usepackage[french]{babel}

% ── Mise en page ──────────────────────────────────────────────
\usepackage[top=2.5cm,bottom=2.5cm,left=3cm,right=3cm]{geometry}
\usepackage{setspace}
\onehalfspacing
\setlength{\parskip}{6pt}
\setlength{\parindent}{0pt}

% ── Typographie ───────────────────────────────────────────────
\usepackage{lmodern}
\usepackage{microtype}

% ── Couleurs ─────────────────────────────────────────────────
\usepackage[table]{xcolor}
\definecolor{primary}{RGB}{30, 90, 168}
\definecolor{rowalt}{RGB}{235, 242, 252}

% ── Tableaux ─────────────────────────────────────────────────
\usepackage{booktabs}
\usepackage{tabularx}
\usepackage{array}
\newcolumntype{L}{>{\raggedright\arraybackslash}X}

% ── Listes ───────────────────────────────────────────────────
\usepackage{enumitem}
\setlist[itemize]{noitemsep, topsep=3pt, leftmargin=1.4em}
\setlist[enumerate]{noitemsep, topsep=3pt, leftmargin=1.6em}

% ── Titres ───────────────────────────────────────────────────
\usepackage{titlesec}
\titleformat{\section}{\Large\bfseries\color{primary}}{}{0em}{}[\vspace{-4pt}\rule{\linewidth}{1.2pt}\vspace{2pt}]
\titleformat{\subsection}{\large\bfseries\color{primary!75!black}}{}{0em}{}
\titleformat{\subsubsection}{\normalsize\bfseries}{}{0em}{}

% ── En-têtes / pieds de page ─────────────────────────────────
\usepackage{fancyhdr}
\pagestyle{fancy}
\fancyhf{}
\fancyhead[L]{\small\textit{Cahier des Charges — SAP Performance Dashboard}}
\fancyhead[R]{\small\thepage}
\fancyfoot[C]{\small\textit{Confidentiel — Usage interne}}
\renewcommand{\headrulewidth}{0.4pt}
\renewcommand{\footrulewidth}{0.4pt}

% ── Boîtes ───────────────────────────────────────────────────
\usepackage{tcolorbox}
\tcbuselibrary{skins, breakable}

\newtcolorbox{rolebox}[2][]{
  colback=primary!6, colframe=primary!70,
  fonttitle=\bfseries\small, title={#2},
  left=6pt, right=6pt, top=4pt, bottom=4pt,
  breakable, arc=4pt, #1
}
\newtcolorbox{notebox}{
  colback=orange!8, colframe=orange!60,
  left=6pt, right=6pt, top=4pt, bottom=4pt,
  arc=4pt, breakable
}

% ── Liens PDF ────────────────────────────────────────────────
\usepackage[hidelinks,
  pdftitle={Cahier des Charges -- SAP Performance Dashboard},
  pdfsubject={Specification fonctionnelle}]{hyperref}

% ─────────────────────────────────────────────────────────────
\begin{document}

% ════════════════════════════════════════════════════════════
\begin{titlepage}
  \centering
  \vspace*{2cm}
  {\color{primary}\rule{\linewidth}{2pt}}\\[0.6em]
  {\Huge\bfseries\color{primary} Cahier des Charges}\\[0.4em]
  {\huge SAP Performance Dashboard}\\[0.4em]
  {\color{primary}\rule{\linewidth}{2pt}}

  \vspace{1.5cm}
  \begin{tcolorbox}[colback=primary!7, colframe=primary, arc=6pt,
                    boxrule=1pt, width=0.80\linewidth, halign=center]
    \large
    Plateforme collaborative de gestion des projets SAP,\\
    des tickets, du temps de travail et de la performance\\
    des equipes de consultants.
  \end{tcolorbox}

  \vspace{2cm}
  \begin{tabular}{@{}ll@{}}
    \textbf{Version :}   & 1.0 \\[4pt]
    \textbf{Date :}      & \today \\[4pt]
    \textbf{Statut :}    & Document de reference \\[4pt]
    \textbf{Auteur(s) :} & Equipe Projet \\
  \end{tabular}

  \vfill
  {\small\color{gray}\textit{Ce document est la propriete exclusive
  de l organisation. Toute reproduction sans autorisation est interdite.}}
\end{titlepage}

\tableofcontents
\newpage

% ════════════════════════════════════════════════════════════
\section{Presentation generale}
% ════════════════════════════════════════════════════════════

\subsection{Contexte}

Les equipes de consultants SAP travaillent simultanement sur plusieurs projets :
projets de maintenance (TMA) et projets de developpement (BUILD). Sans outil
centralise, le suivi des demandes, la declaration des heures travaillees,
la validation des livrables et le pilotage de la performance reposent sur
des echanges d'e-mails et des tableurs eparpilles.

Le \textbf{SAP Performance Dashboard} est une plateforme web qui regroupe
en un seul endroit toutes ces activites, du ticket de developpement jusqu'a
l'export des heures vers les outils RH officiels.

\subsection{Objectifs de la plateforme}

\begin{itemize}
  \item Centraliser la gestion des tickets, des projets et des equipes.
  \item Offrir a chaque acteur une interface adaptee a son role.
  \item Assurer la tracabilite complete de chaque action (qui a fait quoi et quand).
  \item Fluidifier les processus de validation (tickets, imputations, conges, WRICEF).
  \item Donner aux managers et chefs de projet les indicateurs necessaires a la
        prise de decision.
  \item Connecter la plateforme aux outils existants : export vers StraTIME,
        chat Microsoft Teams.
\end{itemize}

\subsection{Utilisateurs de la plateforme}

Six profils d'utilisateurs interagissent avec la plateforme. Chacun dispose
d'une interface dediee, de menus adaptes et d'actions reservees a son role.

\begin{center}
\begin{tabularx}{\linewidth}{|l|L|}
  \hline
  \rowcolor{primary!20}
  \textbf{Role} & \textbf{Responsabilite principale} \\
  \hline
  Administrateur         & Gere les comptes utilisateurs et les parametres globaux \\
  \hline
  Manager                & Pilote les equipes, valide tickets, imputations et conges \\
  \hline
  Chef de projet (PM)    & Supervise les projets, valide les WRICEF et les imputations \\
  \hline
  Coordinateur dev       & Affecte les tickets aux developpeurs, suit les charges \\
  \hline
  Consultant technique   & Realise les developpements, declare son temps de travail \\
  \hline
  Consultant fonctionnel & Cree les tickets de demande, gere les livrables \\
  \hline
\end{tabularx}
\end{center}

% ════════════════════════════════════════════════════════════
\section{Ce que fait chaque utilisateur}
% ════════════════════════════════════════════════════════════

\subsection{Administrateur}

L'administrateur gere les fondations de la plateforme. Il n'intervient pas
dans les projets ou les tickets mais s'assure que les donnees de base sont
correctes.

\begin{rolebox}{L'administrateur peut :}
\begin{itemize}
  \item Creer, modifier ou desactiver des comptes utilisateurs.
  \item Attribuer un role a chaque utilisateur.
  \item Gerer les competences et certifications des consultants.
  \item Administrer les referentiels (statuts, priorites, types de modules SAP).
  \item Consulter un tableau de bord global avec le nombre d'utilisateurs,
        de projets et de tickets en cours.
\end{itemize}
\end{rolebox}

\subsection{Manager}

Le manager a la vue la plus large dans la plateforme.
Il supervise plusieurs equipes et projets en parallele.

\begin{rolebox}{Le manager peut :}
\begin{itemize}
  \item Creer et gerer des projets (TMA ou BUILD) avec leurs dates, priorites
        et equipes affectees.
  \item Consulter un tableau de bord global : productivite, taux d'allocation,
        tickets bloques ou critiques.
  \item Approuver ou rejeter les tickets crees par les consultants fonctionnels,
        en assignant un consultant technique et un volume d'heures.
  \item Valider ou rejeter les declarations de temps (imputations) de son equipe.
  \item Approuver ou refuser les demandes de conge.
  \item Consulter le registre des certifications de chaque consultant.
  \item Importer et gerer les fichiers WRICEF (specifications de developpements SAP).
  \item Affecter les consultants a des projets avec un pourcentage d'allocation.
  \item Voir les risques en temps reel : tickets bloques, en retard ou critiques.
  \item Exporter les imputations validees vers StraTIME (outil RH officiel).
\end{itemize}
\end{rolebox}

\subsection{Chef de projet (PM)}

Le chef de projet pilote un ou plusieurs projets au quotidien.

\begin{rolebox}{Le chef de projet peut :}
\begin{itemize}
  \item Consulter et mettre a jour ses projets (avancement, equipe, documentation).
  \item Voir tous les tickets lies a ses projets et les affecter.
  \item Valider ou rejeter les fichiers WRICEF soumis par le manager.
  \item Valider les imputations de l'equipe projet.
  \item Exporter les imputations vers StraTIME.
  \item Creer et consulter la documentation du projet (SFD, guides, architectures).
  \item Suivre les allocations des consultants sur ses projets.
\end{itemize}
\end{rolebox}

\subsection{Coordinateur developpement}

Le coordinateur est l'intermediaire entre les demandes metier et l'equipe
technique. Il optimise la repartition du travail entre les developpeurs.

\begin{rolebox}{Le coordinateur peut :}
\begin{itemize}
  \item Visualiser tous les tickets en attente d'affectation.
  \item Utiliser l'outil d'aide a l'affectation (IA) : classe les developpeurs
        automatiquement selon charge, competences et historique.
  \item Affecter manuellement ou confirmer l'affectation suggeree par l'IA.
  \item Visualiser la charge de travail actuelle de chaque consultant.
  \item Gerer ses propres declarations d'heures.
\end{itemize}
\end{rolebox}

\subsection{Consultant technique}

Le consultant technique est le developpeur SAP. Il recoit des tickets,
effectue les developpements et declare son temps de travail.

\begin{rolebox}{Le consultant technique peut :}
\begin{itemize}
  \item Consulter ses tickets assignes et mettre a jour leur avancement
        (En cours, En test, Bloque, Termine).
  \item Voir les projets sur lesquels il est affecte.
  \item Saisir ses heures de travail quotidiennes par ticket via un calendrier.
  \item Soumettre ses declarations de temps bi-hebdomadaires pour validation.
  \item Gerer ses demandes de conge et suivre leur statut.
  \item Mettre a jour son profil, ses competences et ses certifications.
  \item Consulter et lier des documents de reference a ses tickets.
\end{itemize}
\end{rolebox}

\subsection{Consultant fonctionnel}

Le consultant fonctionnel represente le metier cote client. Il exprime les
besoins sous forme de tickets et suit l'avancement des livraisons.

\begin{rolebox}{Le consultant fonctionnel peut :}
\begin{itemize}
  \item Creer des tickets de demande (developpement, correction, rapport,
        amelioration) en precisant la priorite, le module SAP et la complexite.
  \item Suivre le statut de ses tickets (en attente d'approbation, en cours, termine).
  \item Consulter les projets sur lesquels il intervient.
  \item Deposer et consulter des livrables (documents, fichiers, liens).
  \item Acceder a la documentation de projet (SFD, guides).
\end{itemize}
\end{rolebox}

% ════════════════════════════════════════════════════════════
\section{Fonctionnalites principales}
% ════════════════════════════════════════════════════════════

\subsection{Gestion des projets}

Chaque projet possede un nom, un type (TMA ou BUILD), un manager, des dates,
une priorite et un pourcentage d'avancement. La fiche projet est organisee
en onglets :

\begin{center}
\begin{tabularx}{\linewidth}{|l|L|}
  \hline
  \rowcolor{primary!20}
  \textbf{Onglet} & \textbf{Contenu} \\
  \hline
  Apercu         & Informations generales, indicateurs cles (tickets, heures, risques) \\
  \hline
  Objets (WRICEF) & Liste des specifications de developpements liees \\
  \hline
  Tickets        & Tous les tickets du projet avec filtres et actions rapides \\
  \hline
  Equipe         & Consultants affectes avec leur pourcentage d'allocation \\
  \hline
  KPI            & Graphiques de suivi : progression, charge, statuts \\
  \hline
  Documentation  & Base de connaissances du projet (SFD, guides, architectures) \\
  \hline
  Abaques        & Grille d'estimation par nature et complexite de ticket \\
  \hline
\end{tabularx}
\end{center}

\subsection{Gestion des tickets}

\subsubsection{Qu'est-ce qu'un ticket ?}

Un ticket represente une demande de travail : developpement d'un programme SAP,
correction d'un bug, creation d'un formulaire, generation d'un rapport, etc.
Chaque ticket recoit un code unique automatique (exemple : TK-2026-AB3F9C).

\subsubsection{Cycle de vie d'un ticket}

\begin{center}
\begin{tabularx}{\linewidth}{|l|l|L|}
  \hline
  \rowcolor{primary!20}
  \textbf{Statut} & \textbf{Declencheur} & \textbf{Signification} \\
  \hline
  En attente d'approbation & Creation par le consultant fonctionnel
    & Le manager doit approuver la demande \\
  \hline
  Nouveau        & Approbation du manager   & Pret a etre pris en charge \\
  \hline
  En cours       & Prise en charge          & Developpement actif \\
  \hline
  En test        & Developpement termine    & En attente de recette fonctionnelle \\
  \hline
  Termine        & Validation finale        & Livraison confirmee \\
  \hline
  Bloque         & Obstacle identifie       & Intervention necessaire \\
  \hline
  Rejete         & Refus par le manager     & Demande non acceptee (avec motif) \\
  \hline
\end{tabularx}
\end{center}

\subsubsection{Informations d'un ticket}

\begin{itemize}
  \item Projet, titre, description, priorite, module SAP, complexite, nature.
  \item Estimation des heures (saisie manuelle ou via la grille Abaque).
  \item Consultant assigne et testeur fonctionnel.
  \item Lien vers la specification WRICEF correspondante.
  \item Historique horodate de toutes les modifications.
  \item Fil de commentaires et documents lies.
\end{itemize}

\subsubsection{Filtres disponibles}

Les listes de tickets peuvent etre filtrees par statut, module SAP, complexite,
priorite, projet, consultant assigne, reference WRICEF, plage de dates et
recherche textuelle.

\subsection{Commentaires et collaboration}

Chaque ticket dispose d'un fil de discussion structure. Un commentaire peut etre
de type : General, Bloquant, Question, Mise a jour ou Retour. Les commentaires
peuvent etre replies en fil groupe et marques comme resolus. Il est aussi possible
d'ouvrir directement une conversation Microsoft Teams avec le responsable du
ticket sans quitter la plateforme.

\subsection{Gestion des WRICEF}

Les WRICEF (Workflow, Report, Interface, Conversion, Enhancement, Form) sont les
specifications formelles des developpements SAP. Chaque ligne WRICEF decrit un
objet a developper avec son type, sa complexite et son module SAP.

\begin{enumerate}
  \item Le manager importe un fichier Excel ou CSV contenant la liste des objets WRICEF.
  \item L'application analyse le fichier et presente un apercu avant validation.
  \item Une fois importes, les objets WRICEF aparaissent dans la fiche projet.
  \item Le manager soumet le dossier pour validation au chef de projet.
  \item Le chef de projet valide ou rejette le dossier (avec motif si rejet).
\end{enumerate}

Etats d'un dossier WRICEF : \textbf{Brouillon} \textrightarrow{}
\textbf{Soumis} \textrightarrow{} \textbf{Valide / Rejete}.

\subsection{Declaration des heures (Imputations)}

Chaque consultant technique declare ses heures de travail au quotidien, par ticket.
Ces declarations sont regroupees automatiquement par quinzaines (1\textsuperscript{re}
et 2\textsuperscript{e} moitie du mois) et soumises a validation.

\begin{enumerate}
  \item Le consultant saisit ses heures chaque jour via un calendrier mensuel.
  \item En fin de quinzaine, il soumet sa periode pour approbation.
  \item Le manager ou chef de projet valide ou renvoie les declarations en correction.
  \item Une fois validees, les imputations sont exportees vers \textbf{StraTIME}.
\end{enumerate}

\begin{notebox}
  Une declaration soumise ou validee ne peut plus etre supprimee.
  Seules les declarations en brouillon restent modifiables.
\end{notebox}

\subsection{Gestion documentaire}

Chaque projet dispose d'une base de connaissances accessible a toute l'equipe.
Les documents peuvent etre de quatre types :

\begin{itemize}
  \item \textbf{SFD} — Specification Fonctionnelle Detaillee d'un developpement.
  \item \textbf{Guide utilisateur} — documentation d'usage d'un outil ou processus.
  \item \textbf{Document d'architecture} — vision d'ensemble du systeme.
  \item \textbf{General} — tout autre document utile a l'equipe.
\end{itemize}

Chaque document peut etre lie a des tickets et a des WRICEF, et recevoir des
pieces jointes (fichiers, liens).

\subsection{Livrables}

Un livrable est un document produit dans le cadre d'un projet (rapport de recette,
guide, schema, etc.). Le consultant fonctionnel depose son livrable sur la
plateforme, qui suit son statut de revue :
\textbf{En attente} \textrightarrow{} \textbf{Approuve / A corriger}.

\subsection{Allocations des ressources}

Le manager affecte chaque consultant a un ou plusieurs projets en precisant
le pourcentage de temps alloue (ex. 50\% sur le projet A, 50\% sur le projet B)
et la periode de mission. Cette information est visible sur la fiche projet et
utilisee par l'outil d'aide a l'affectation IA.

\subsection{Conges}

\begin{enumerate}
  \item Le consultant soumet une demande de conge (dates et motif).
  \item Le manager approuve ou refuse la demande.
  \item Le consultant est notifie de la decision.
\end{enumerate}

\subsection{Evaluations}

Le manager ou le chef de projet peut evaluer un consultant a la fin d'un projet
ou d'une periode, avec un score global, une grille qualitative par critere
(qualite, delais, communication...) et un commentaire libre.

\subsection{Aide a l'affectation par intelligence artificielle}

Le coordinateur developpement dispose d'un outil d'aide a la decision pour
affecter les tickets non assignes. L'outil analyse automatiquement les consultants
disponibles selon quatre criteres :

\begin{center}
\begin{tabularx}{\linewidth}{|l|c|L|}
  \hline
  \rowcolor{primary!20}
  \textbf{Critere} & \textbf{Poids} & \textbf{Description} \\
  \hline
  Disponibilite          & 30\% & Taux d'occupation actuel du consultant \\
  \hline
  Competences            & 35\% & Adequation avec la nature du ticket \\
  \hline
  Performance historique & 15\% & Score moyen des evaluations passees \\
  \hline
  Tickets similaires     & 20\% & Nombre de tickets du meme type deja traites \\
  \hline
\end{tabularx}
\end{center}

Le resultat est un classement des consultants avec une explication lisible
pour chaque candidat. Le coordinateur confirme ou modifie le choix propose.

\subsection{Abaques d'estimation}

Un abaque est une grille d'estimation standard definie par projet. Pour chaque
combinaison de \textit{nature de ticket} et de \textit{complexite}, l'abaque
donne un nombre d'heures de reference. Lors de la creation d'un ticket,
l'utilisateur peut choisir d'utiliser l'abaque : les heures sont alors
pre-remplies automatiquement.

\subsection{Tableaux de bord}

Chaque role beneficie d'un tableau de bord adapte a ses besoins :

\begin{center}
\begin{tabularx}{\linewidth}{|l|L|}
  \hline
  \rowcolor{primary!20}
  \textbf{Role} & \textbf{Indicateurs affiches} \\
  \hline
  Administrateur &
    Nombre d'utilisateurs actifs, projets en cours, tickets ouverts \\
  \hline
  Manager &
    Productivite de l'equipe, taux d'allocation, tickets bloques ou critiques,
    heures declarees vs. estimees, demandes en attente \\
  \hline
  Chef de projet &
    Avancement des projets pilotes, tickets par statut, charge de l'equipe \\
  \hline
  Coordinateur dev &
    Tickets non assignes, charge par consultant, repartition par module \\
  \hline
  Consultant technique &
    Mes tickets actifs, heures declarees ce mois, avancement personnel \\
  \hline
  Consultant fonctionnel &
    Mes tickets en cours, statut des livrables, avancement des projets \\
  \hline
\end{tabularx}
\end{center}

\subsection{Notifications}

La plateforme envoie des notifications internes lors d'evenements cles
(approbation d'un ticket, validation d'une imputation, decision sur un conge).
Les notifications apparaissent dans la barre superieure avec un compteur
de messages non lus.

% ════════════════════════════════════════════════════════════
\section{Processus metier detailles}
% ════════════════════════════════════════════════════════════

\subsection{Processus 1 — Creation et approbation d'un ticket}

\begin{enumerate}
  \item Le \textbf{consultant fonctionnel} remplit le formulaire : titre,
        description, priorite, module SAP, complexite, lien WRICEF.
  \item Le ticket est cree avec le statut \textit{En attente d'approbation}.
  \item Le \textbf{manager} recoit une notification et accede a la page
        \textit{Approbations en attente}.
  \item \textbf{Si approuve :} le manager choisit un consultant technique et
        les heures allouees. Le ticket passe au statut \textit{Nouveau}.
  \item \textbf{Si rejete :} le manager saisit un motif. Le consultant
        fonctionnel est notifie.
\end{enumerate}

\subsection{Processus 2 — Import et validation d'un WRICEF}

\begin{enumerate}
  \item Le \textbf{manager} importe un fichier Excel/CSV depuis la fiche projet.
  \item L'application presente la liste des objets detectes avec un apercu.
  \item Le manager confirme l'import \textrightarrow{} objets enregistres en
        statut \textit{Brouillon}.
  \item Le manager soumet le dossier \textrightarrow{} statut \textit{En attente
        de validation}.
  \item Le \textbf{chef de projet} valide ou rejette le dossier global.
\end{enumerate}

\subsection{Processus 3 — Declaration et validation des imputations}

\begin{enumerate}
  \item Le \textbf{consultant technique} saisit ses heures chaque jour :
        ticket, duree, description.
  \item En fin de quinzaine, il soumet sa periode.
  \item Le \textbf{manager} ou le \textbf{chef de projet} valide ou renvoie
        en correction.
  \item Une fois validees, les imputations sont exportees vers \textbf{StraTIME}.
\end{enumerate}

\subsection{Processus 4 — Affectation assistee d'un ticket (IA)}

\begin{enumerate}
  \item Le \textbf{coordinateur} consulte la liste des tickets non assignes.
  \item Il selectionne un ticket et lance l'analyse.
  \item L'outil presente un classement des consultants avec le detail du score.
  \item Le coordinateur confirme ou modifie le choix, puis valide l'affectation.
\end{enumerate}

% ════════════════════════════════════════════════════════════
\section{Integrations avec les outils externes}
% ════════════════════════════════════════════════════════════

\subsection{StraTIME — Export des heures}

StraTIME est le systeme officiel de declaration des temps. La plateforme permet
d'exporter les imputations validees directement vers cet outil d'un simple clic,
evitant toute double saisie. Chaque imputation exportee est marquee pour eviter
les doublons en cas de renvoi.

\subsection{Microsoft Teams — Chat integre}

Depuis la page d'un ticket, tout utilisateur peut ouvrir directement une
conversation Microsoft Teams avec le responsable ou le testeur du ticket.
Le message est pre-rempli avec le contexte du ticket, facilitant la communication
sans quitter la plateforme.

% ════════════════════════════════════════════════════════════
\section{Regles metier importantes}
% ════════════════════════════════════════════════════════════

\begin{itemize}
  \item Un consultant ne voit que ses propres tickets.
  \item Seul un manager peut approuver ou rejeter un ticket.
  \item Un dossier WRICEF doit contenir au moins un objet avant d'etre soumis.
  \item Seul le chef de projet peut valider un dossier WRICEF.
  \item Une imputation soumise ou validee ne peut plus etre supprimee.
  \item L'export vers StraTIME ne concerne que les periodes dont le statut
        est \textit{Valide}.
  \item Les roles Manager et Administrateur ne sont jamais proposes
        par l'outil d'affectation IA.
\end{itemize}

% ════════════════════════════════════════════════════════════
\section*{Conclusion}
\addcontentsline{toc}{section}{Conclusion}
% ════════════════════════════════════════════════════════════

Le \textbf{SAP Performance Dashboard} couvre l'integralite du cycle de vie
d'un projet SAP : de la specification initiale (WRICEF) jusqu'a la livraison
du ticket, en passant par la declaration des heures, les evaluations et
l'export vers les systemes RH. Chaque acteur dispose d'un espace de travail
adapte a ses responsabilites, et tous les processus de validation sont
tracables et auditables.

Ce document constitue le referentiel fonctionnel de base pour la conception
detaillee, la recette et la conduite du changement aupres des utilisateurs.

\end{document}
